% Source: https://artofproblemsolving.com/community/c3h2796071_2012_tx_mmc__texas_mathworks_math_contest
% Stable identity: aops-topic-2796071-q09
Let $AB=11$ units, so $EB=2$, $EA=9$, $DA=DB=5.5$ and $DE=9-5.5=3.5$. Let $h$ be the height of $C$ above $AB$.
\begin{asy}
usepackage("amsmath");
size(5.5cm);
defaultpen(fontsize(10pt));
pair A=(0,0),B=(11,0),C=(3.5,6);
pair Dm=(A+B)/2,Ep=(9,0),F=B+(2/7)*(C-B);
fill(Dm--Ep--F--cycle,gray(0.8));
draw(A--B--C--cycle);draw(Dm--F--Ep);draw(A--B);draw(Dm--C,dashed);
label("$A$",A,SW);label("$B$",B,SE);label("$C$",C,N);label("$D$",Dm,S);label("$E$",Ep,S);label("$F$",F,NE);
\end{asy}

Since $BF=\dfrac27BC$, the height of $F$ above $AB$ is $\dfrac27h$. Thus
\[[DEF]=\dfrac12\cdot3.5\cdot\dfrac27h=\dfrac h2=\dfrac1{11}\cdot\dfrac{11h}{2}=\dfrac{[ABC]}{11}.\]
So $[ABC]=132$, and since $D$ is the midpoint of $AB$, $[ADC]=\dfrac{132}2=\boxed{66}$.
